%% titre: TD 1 — Espaces vectoriels
%% niveau: Licence 1
%% matiere: Algèbre linéaire
%% annee: 2025-2026
%% ordre: 1
%% description: Exemple de feuille de TD (à remplacer) avec exercices et corrections.

\begin{exercice}
Dans $\R^3$, les ensembles suivants sont-ils des sous-espaces vectoriels ?
\begin{enumerate}
  \item $E_1=\{(x,y,z)\in\R^3 \mid x+y+z=0\}$ ;
  \item $E_2=\{(x,y,z)\in\R^3 \mid x+y+z=1\}$ ;
  \item $E_3=\{(x,y,z)\in\R^3 \mid xy=0\}$.
\end{enumerate}
\begin{correction}
\begin{enumerate}
  \item Oui : c'est le noyau de la forme linéaire $(x,y,z)\mapsto x+y+z$.
  \item Non : le vecteur nul n'appartient pas à $E_2$.
  \item Non : $(1,0,0)$ et $(0,1,0)$ sont dans $E_3$ mais leur somme $(1,1,0)$ n'y est pas.
\end{enumerate}
\end{correction}
\end{exercice}

\begin{exercice}[$\bigstar$]
Soit $f:\R^2\to\R^2$ définie par $f(x,y)=(x+2y,\,2x+4y)$.
\begin{enumerate}[label=\alph*)]
  \item Montrer que $f$ est linéaire.
  \item Déterminer $\Ker f$ et $\Im f$.
  \item $f$ est-elle injective ? surjective ?
\end{enumerate}
\begin{correction}
La matrice de $f$ dans la base canonique est $A=\begin{pmatrix}1&2\\2&4\end{pmatrix}$, de rang $1$.
Donc $\Ker f=\Vect\{(-2,1)\}$ et $\Im f=\Vect\{(1,2)\}$. Comme $\Ker f\neq\{0\}$, $f$ n'est pas injective ; comme $\rg f=1<2$, elle n'est pas surjective.
\end{correction}
\end{exercice}

\begin{exercice}
Soit $A\in\mathcal{M}_n(\R)$ telle que $A^2=A$. Montrer que
\[
  \R^n=\Ker A\oplus\Im A .
\]
\begin{correction}
Pour $x\in\R^n$, on écrit $x=(x-Ax)+Ax$ avec $x-Ax\in\Ker A$ car $A(x-Ax)=Ax-A^2x=0$. Si $y\in\Ker A\cap\Im A$, alors $y=Az$ et $0=Ay=A^2z=Az=y$. La somme est donc directe.
\end{correction}
\end{exercice}
